\begin{pastexam}[label={px:df3-100-ds-01}]{地方三等資訊處理 100 年 第一題}
（申論；20 分）

給一規則的整數數列 A=\{A$_0$, A$_1$, A$_2$, ...\}=\{1, 1, 1, 3, 5, 9, 17, 31, 57,…\}。
\begin{itemize}[leftmargin=2.2em, nosep]
\item[㈠] 試寫一遞迴函式（recursive function）計算 A$_n$ 的數值。（10 分）
\item[㈡] 利用上述遞迴方法詳列計算 A$_6$ 數值的過程。（10 分）
\end{itemize}
\end{pastexam}

\begin{pastsoln}
答案：$A_n = A_{n-1} + A_{n-2} + A_{n-3}$（$n \ge 3$），$A_0 = A_1 = A_2 = 1$；$A_6 = 17$。

\textbf{找規則。} 從第四項起，每一項是前三項的和：$3 = 1 + 1 + 1$，$5 = 1 + 1 + 3$，$9 = 1 + 3 + 5$，
$17 = 3 + 5 + 9$，$31 = 5 + 9 + 17$，$57 = 9 + 17 + 31$，全部吻合。 所以
\[
A_n = \begin{cases} 1, & n = 0, 1, 2, \\ A_{n-1} + A_{n-2} + A_{n-3}, & n \ge 3. \end{cases}
\]

\textbf{㈠} 照定義直接寫成遞迴，$n \le 2$ 是終止條件：
\begin{minted}[fontsize=\small]{c}
long long A(int n)
{
    if (n <= 2)          /* A0 = A1 = A2 = 1 */
        return 1;
    return A(n - 1) + A(n - 2) + A(n - 3);
}
\end{minted}

\textbf{㈡} 呼叫 $A(6)$ 時，函式先依序呼叫 $A(5)$、$A(4)$、$A(3)$，每一個再往下展開，直到參數 $\le 2$
傳回 $1$。 展開過程（由上往下呼叫，由下往上傳回）：
\begin{align*}
A(6) &= A(5) + A(4) + A(3) \\
A(5) &= A(4) + A(3) + A(2) \\
A(4) &= A(3) + A(2) + A(1) \\
A(3) &= A(2) + A(1) + A(0) = 1 + 1 + 1 = 3
\end{align*}
再一層一層代回去：
\begin{align*}
A(4) &= A(3) + A(2) + A(1) = 3 + 1 + 1 = 5 \\
A(5) &= A(4) + A(3) + A(2) = 5 + 3 + 1 = 9 \\
A(6) &= A(5) + A(4) + A(3) = 9 + 5 + 3 = 17
\end{align*}
所以 $A_6 = 17$。

要注意的是遞迴版並不會「記住」算過的值：$A(6)$ 呼叫的 $A(4)$ 和 $A(5)$ 裡面呼叫的 $A(4)$ 是兩次獨立的
計算。 各參數被呼叫的次數：
\begin{center}
\begin{tabular}{lcccccccc}
\toprule
參數 & 6 & 5 & 4 & 3 & 2 & 1 & 0 & 合計 \\
\midrule
被呼叫次數 & 1 & 1 & 2 & 4 & 7 & 6 & 4 & 25 \\
\bottomrule
\end{tabular}
\end{center}
（$A(4)$ 被 $A(6)$、$A(5)$ 各呼叫一次；$A(3)$ 被 $A(6)$、$A(5)$ 各一次，再被兩次 $A(4)$ 各一次，共 $4$ 次；
以此類推。） 呼叫次數 $C(n) = 1 + C(n-1) + C(n-2) + C(n-3)$ 和 $A_n$ 一樣呈指數成長（資結書第~2~章），
所以 $n$ 大時要改成由小到大用迴圈算，或把算過的值存起來（資結書第~18.1~節），時間變成 $O(n)$。
\end{pastsoln}
